\section{Related Work}
\label{sec:related_work}

Network intrusion detection has evolved from manually specified signatures to
data-driven models that learn representations from packet traces, flow records,
or sequences of network events. The literature can be organized along two
closely related dimensions: the granularity of the traffic representation and
the scope of the modeled context. In terms of granularity, existing systems use
packet bytes, packet-header sequences, aggregated flow statistics, or learned
flow embeddings. In terms of scope, most detectors classify each packet or flow
independently, whereas a smaller body of work models temporal or relational
dependencies among multiple flows. This section reviews these directions and
positions the proposed hierarchical two-stage model with respect to them.

\subsection{Signature-Based and Statistical Flow-Based Detection}
\label{subsec:rw_signature_flow}

Signature-based NIDSs identify traffic patterns that match expert-defined
rules. Snort established a lightweight and extensible architecture for
real-time packet inspection \cite{Roesch1999}, while Suricata provides a modern
multi-threaded engine and programmable logging facilities
\cite{SuricataGuide2026}. These systems remain important in operational
networks because their alerts are interpretable and their behavior is
controllable. However, their effectiveness depends on the availability and
maintenance of signatures, and previously unseen or behaviorally modified
attacks may not match known rules. More generally, Sommer and Paxson caution
that learning-based NIDS research must account for the open-world character of
network traffic and the semantic gap between laboratory benchmarks and
operational deployment \cite{SommerPaxson2010}.

Classical machine-learning NIDSs commonly transform packets into bidirectional
flow records and represent each flow by statistics such as duration, packet and
byte counts, rates, TCP flags, and inter-arrival-time summaries. Surveys by
Ahmed et al. and Buczak and Guven describe the broad use of decision trees,
support-vector machines, nearest-neighbor methods, clustering, and ensemble
models for anomaly and intrusion detection
\cite{Ahmed2016Survey,BuczakGuven2016}. Feature selection and hybrid
classification can reduce dimensionality and improve conventional models
\cite{Aljawarneh2018}. The principal advantage of flow aggregation is
efficiency: one fixed-dimensional vector represents a variable-length network
conversation. Its principal limitation is information loss. Two flows with
similar totals may exhibit different packet orders, direction changes, burst
structures, or timing gaps. Consequently, a detector trained only on aggregate
statistics cannot directly learn the packet-level dynamics that generated
those statistics.

The present work retains flow statistics because they provide a compact summary
of volume and protocol behavior, but it does not treat them as a sufficient
representation. Instead, they are used as a complementary modality and fused
with an embedding learned from the packet sequence. This design also motivates
the flow-statistics MLP baseline used to test whether packet-level temporal
modeling contributes information beyond aggregation alone.

\subsection{Deep Learning for Packet and Flow Sequences}
\label{subsec:rw_deep_sequences}

Deep neural models reduce dependence on manually selected features by learning
non-linear representations of network traffic. Recurrent NIDSs use hidden state
to summarize ordered observations; for example, Yin et al. applied recurrent
neural networks to intrusion detection and demonstrated their ability to learn
sequential dependencies \cite{Yin2017RNNIDS}. Shone et al. combined a
non-symmetric deep autoencoder with supervised classification to learn compact
traffic representations \cite{Shone2018NNAE}. Kitsune instead used an ensemble
of autoencoders for lightweight online anomaly detection, emphasizing the
deployment value of compact and incrementally processed features
\cite{Mirsky2018Kitsune}. Hybrid CNN--LSTM designs have also been used to combine
local pattern extraction with recurrent temporal modeling
\cite{Ullah2024IDSINT}.

Sequence construction is as important as the choice of neural architecture.
Many studies apply an LSTM or GRU to a vector of flow features, where the
``sequence'' corresponds to feature dimensions rather than chronologically
ordered packets or flows. Such a formulation can learn feature interactions but
does not necessarily model network time. NetSentry addresses this issue more
directly by treating intrusion detection as a time-sensitive task and using
ordered network observations to detect early-stage large-scale attacks
\cite{LiuPatras2022NetSentry}. Its results support the broader argument that
temporally correlated behavior can be more informative than independently
classified records. Nevertheless, recurrence imposes sequential computation,
and a single final hidden state may become a bottleneck for long or irregularly
sampled sequences.

The Stage1 baselines in this thesis therefore include CNN, LSTM, GRU, BiLSTM,
and time-augmented recurrent variants. These comparisons separate the value of
the proposed traffic representation from the value of the Transformer
architecture itself. In particular, they test whether explicit packet timing
and attention-based aggregation improve minority-class detection relative to
standard sequential encoders.

\subsection{Transformer-Based Traffic Representation}
\label{subsec:rw_transformers}

Transformers replace recurrent state transitions with self-attention, enabling
each token to interact directly with all other tokens in a sequence and
supporting parallel training \cite{Vaswani2017}. This mechanism is attractive
for network traffic because attack evidence may be separated by many packets
or flows. RTIDS demonstrated that a Transformer-based detector can provide
robust intrusion classification on flow-oriented benchmarks
\cite{Wu2022RTIDS}. Long et al. subsequently studied an encoder-based
Transformer for cloud network intrusion detection
\cite{Long2024TransformerNIDS}, while IDS-INT combined Transformer-based
transfer learning with imbalance handling and downstream neural components
\cite{Ullah2024IDSINT}.

Two studies are especially close to the representation problem considered in
Stage1. Han et al. proposed GTID, which combines packet- and session-level
features with a time-aware Transformer; importantly, its temporal encoding uses
inter-packet intervals rather than relying only on ordinal position
\cite{Han2023GTID}. FlowTransformer provides a modular framework for examining
input encodings, Transformer blocks, classification heads, model size, and
runtime on flow-based NIDS datasets \cite{Manocchio2024FlowTransformer}. More
recently, the memory flow Transformer was proposed to retain longer-term flow
information and model complex relations in large sequential inputs
\cite{Jiang2025MFT}. Collectively, these works establish that attention-based
models can capture dependencies that are difficult to represent with a single
aggregate vector. They also show that input encoding, pooling, sequence
construction, and efficiency can matter as much as the nominal model family.

Despite this progress, three distinctions remain. First, positional order and
elapsed network time are not equivalent: two adjacent packets may be separated
by microseconds or seconds. Second, packet-sequence features and flow-level
statistics describe complementary aspects of a connection and should be
evaluated both separately and jointly. Third, a representation learned for one
flow does not by itself explain whether that flow forms part of a distributed
host-level pattern. The proposed Stage1 addresses the first two distinctions by
combining positional and continuous-time encodings and by late-fusing the
resulting packet-sequence embedding with flow statistics. Stage2 addresses the
third by contextualizing the learned flow embeddings.

\subsection{Self-Supervised and Encrypted-Traffic Representation Learning}
\label{subsec:rw_ssl_encrypted}

The increasing use of encryption has shifted traffic analysis away from direct
payload inspection toward packet sizes, directions, timing, header metadata,
and learned representations. Recent surveys document this transition and the
growing role of deep pre-training for encrypted traffic classification
\cite{Sharma2025EncryptedSurvey,Dong2025EncryptedPretrain}. ET-BERT adapts
bidirectional Transformer pre-training to contextualized datagram
representations and demonstrates transfer to encrypted-traffic classification
tasks \cite{Lin2022ETBERT}. Masked-autoencoder approaches similarly learn from
unlabeled traffic by reconstructing hidden parts of an input representation
\cite{Zhao2024MAETraffic,2024MAETraffic}. Contrastive learning provides another
route to label-efficient NIDS by bringing related traffic views together while
separating incompatible examples \cite{Li2024Contrastive}.

Other recent systems broaden the representation space. CoTNeT applies a
contextual Transformer mechanism to encrypted-traffic features
\cite{Huang2024CoTNet}, and DoLLM represents flow records in a form suitable for
large-language-model-based DDoS detection \cite{Li2024DoLLM}. These approaches
show that useful traffic representations can be learned without plaintext
payloads and with reduced dependence on manually engineered signatures.
However, encrypted-traffic classification and intrusion detection are not
identical tasks: application identity may be stable even when attack behavior
is sparse, evolving, or distributed across connections. The present thesis is
therefore payload-independent but explicitly supervised for flow-level attack
detection. Self-supervised pre-training is treated as a promising extension
rather than a prerequisite of the proposed two-stage architecture.

\subsection{Temporal, Relational, and Cross-Flow Context}
\label{subsec:rw_cross_flow}

Most flow-based detectors assume that samples are conditionally independent.
This simplification is convenient for batching and online inference, but it is
often inconsistent with attack behavior. Scanning, password guessing, command
and control, DDoS, and lateral movement can produce multiple individually
ambiguous flows whose shared source, destination, endpoint, or temporal pattern
is discriminative. Temporal models such as NetSentry demonstrate the value of
ordered observations \cite{LiuPatras2022NetSentry}, while graph-based NIDSs make
host and flow relations explicit. Deng et al. model flow topology with graph
convolution for label-limited IoT intrusion detection
\cite{Deng2023FlowTopology}; Xu et al. combine graph neural networks with
self-supervised learning for network-flow intrusion detection
\cite{Xu2024SSLGNN}. A recent survey by Zhong et al. organizes the rapidly
growing use of graph neural networks in IDSs and highlights their ability to
encode structural dependencies \cite{Zhong2024GNNSurvey}.

Graph methods provide expressive relational inductive bias, but constructing
and updating a large communication graph can increase memory, latency, and
deployment complexity. Window-based sequence models offer a complementary
design point. A recent Transformer approach aggregates IoT flows into temporal
windows and treats each flow as a token, allowing self-attention to capture
inter-flow dependencies \cite{MartinReizabal2026FlowWindow}. The latest review
of contextualized NetFlow NIDSs classifies context into temporal, relational,
multimodal, and multi-resolution forms and emphasizes that evaluation must
preserve causality and prevent cross-split context leakage
\cite{ElMahdaouy2026ContextNetFlow}.

The Stage2 model in this thesis follows the window-based direction but differs
in two respects. First, its tokens are not raw flow records; they are fused
Stage1 embeddings that already encode packet-level temporal dynamics and flow
statistics. Second, context membership is explicitly controlled through
time-only, source-host, destination-host, and endpoint strategies. This enables
a direct ablation of which relational assumption contributes useful evidence.
The comparison with a no-context classifier tests whether gains originate from
neighboring flows rather than from additional classifier capacity, while the
LSTM, GRU, and CNN--LSTM Stage2 baselines test whether self-attention is needed
to aggregate that context.

Table~\ref{tab:rw_model_comparison} summarizes the modeling scope of several
representative Transformer-based approaches. The comparison is qualitative
rather than performance-based because the studies use different datasets,
labels, preprocessing pipelines, and evaluation protocols. Here,
\emph{explicit elapsed time} means that timestamps or inter-arrival intervals
are encoded directly; positional order or membership in a temporal window alone
does not satisfy this criterion.

\begin{table}[t]
\centering
\caption{Qualitative comparison of representative Transformer-based traffic
models. Entries describe the principal configuration reported in each cited
study; extensible frameworks may support additional variants.}
\label{tab:rw_model_comparison}
\scriptsize
\setlength{\tabcolsep}{2pt}
\renewcommand{\arraystretch}{1.22}
\begin{tabular}{@{}p{1.85cm}p{2.65cm}p{1.55cm}p{2.15cm}p{2.45cm}p{1.85cm}p{1.85cm}@{}}
\toprule
\textbf{Method} &
\textbf{Primary input} &
\textbf{Intra-flow packet order} &
\textbf{Explicit elapsed time} &
\textbf{Inter-flow context} &
\textbf{Payload dependency} &
\textbf{Primary task} \\
\midrule
FlowTransformer \cite{Manocchio2024FlowTransformer} &
Flow records and statistics &
No &
Not explicitly encoded &
Yes; ordered flow sequence &
No &
Flow-level NIDS \\

GTID \cite{Han2023GTID} &
Packet headers and payload n-grams within a session &
Yes &
Yes; inter-packet intervals &
No &
Yes; payload n-grams &
NIDS \\

ET-BERT \cite{Lin2022ETBERT} &
Raw datagram bytes grouped into BURSTs &
Yes; byte and BURST order &
No; order only &
No cross-flow context &
Yes; encrypted bytes &
Encrypted-traffic classification \\

Flow-to-window Transformer \cite{MartinReizabal2026FlowWindow} &
Flow records in fixed-duration windows &
No &
Window-based order, not packet intervals &
Yes; temporal window &
No &
IoT NIDS and traffic classification \\

\textbf{Proposed} &
Packet-header sequence and flow statistics &
\textbf{Yes} &
\textbf{Yes; packet timestamps or intervals} &
\textbf{Yes; time, source, destination, or endpoint} &
\textbf{No direct payload} &
\textbf{Flow-level NIDS} \\
\bottomrule
\end{tabular}
\end{table}

The table exposes the research gap without implying that the compared systems
are directly rankable. FlowTransformer and the flow-to-window model capture
relationships across flow records, but do not reconstruct the ordered packet
dynamics inside each flow. GTID explicitly models packet order and inter-packet
time, but operates on one session at a time and uses payload-derived n-grams.
ET-BERT learns strong byte- and BURST-level representations from encrypted raw
traffic, but its principal objective is encrypted-traffic classification rather
than cross-flow intrusion detection. Among the representative systems in the
table, the proposed framework is the only one that jointly combines
payload-independent intra-flow packet modeling, explicit elapsed-time encoding,
flow-statistical fusion, and configurable inter-flow context. This comparison
therefore motivates the hierarchical separation between Stage1 and Stage2.

\subsection{Datasets, Class Imbalance, and Evaluation Reliability}
\label{subsec:rw_evaluation}

Reported NIDS performance depends strongly on dataset construction and the
evaluation protocol. Ring et al. distinguish packet- and flow-based datasets
and identify recording environment, labeling, traffic diversity, and attack
coverage as central dataset properties \cite{Ring2019Datasets}. More recent
reviews show that the number of public NIDS datasets has grown substantially,
but many datasets still contain limited realism, duplicated records, labeling
errors, outdated attacks, or tool-dependent feature artifacts
\cite{Goldschmidt2025Datasets,Pinto2025DatasetReview}. These problems can make
random record-level splits overly optimistic when highly related flows, hosts,
or capture periods occur in both training and test sets.

Evaluation pitfalls are not limited to the data source. Arp et al. show that
machine-learning security studies are vulnerable to leakage, inappropriate
baselines, unrealistic assumptions, and metrics that obscure the actual
security objective \cite{Arp2022DosDonts}. Class imbalance is particularly
important in NIDS because high accuracy may coexist with poor attack recall.
Focal loss is one possible training mechanism for emphasizing difficult and
minority examples \cite{Lin2017FocalLoss}, but imbalance must also be reflected
in reporting. Precision--recall analysis is generally more informative than
ROC analysis when the positive class is rare
\cite{SaitoRehmsmeier2015}. Accordingly, this thesis treats macro-F1 as the
primary summary metric and reports macro-precision, macro-recall, attack-class
precision/recall/F1, PR-AUC, ROC-AUC, the confusion matrix, and the false
positive rate (FPR). Parameter count, training time, and inference time are also
reported to expose the operational cost of contextual modeling.

To reduce leakage risk, preprocessing parameters and decision thresholds should
be selected without test-set access, and context windows should obey the chosen
split policy. In particular, an online setting may use only information
available before the target flow, whereas split-isolated evaluation must not
allow validation or test flows to retrieve training-incompatible neighbors.
These controls are essential for determining whether Stage2 learns transferable
behavior or merely memorizes hosts, capture periods, or near-duplicate flows.

\subsection{Summary and Research Gap}
\label{subsec:rw_gap}

Prior research has established the usefulness of aggregate flow statistics,
deep sequential encoders, time-aware attention, self-supervised traffic
representations, and graph or window-based context. Nevertheless, these
directions are usually studied separately. Flow-statistical models are
efficient but discard packet dynamics; packet-sequence models usually classify
one flow at a time; and cross-flow models commonly operate on raw flow vectors
without first learning a timing-aware representation inside each flow.

This thesis addresses that gap with a hierarchical division of labor. Stage1
models packet order and elapsed time within each flow and late-fuses the learned
sequence representation with flow statistics. Stage2 then models dependencies
among the resulting flow embeddings under explicit temporal and endpoint-based
context policies. The design therefore connects two levels of network behavior:
\emph{intra-flow} packet dynamics and \emph{inter-flow} attack context. The
ablation studies are organized around the same distinction, allowing the work
to test separately whether time-aware packet modeling improves the flow
representation and whether cross-flow context improves the final detector.
